%% =====================================================================
%%  Capítulo 6 -- Error en estado estacionario
%%  Migrado al formato nuevo (noviembre de 2026) desde Handout5_Ctrl.tex.
%%  Correcciones de fondo: ver MIGRACION.md y la caja "Correcciones" de la
%%  sección 6.8 (fórmulas de kv y ka, retroalimentación no unitaria con el
%%  lazo abierto G_OL, sensitividad complementaria T en lugar de C).
%% =====================================================================

\capitulo[cap:er]{Error en estado estacionario}
  {Los lectores del tercer mundo quieren que alguien les arme una buena ficción con la pésima realidad que les ha tocado en mala suerte.}
  {Rodrigo Fresán}

\begin{objetivos}[partes/06-error.html]
  \item Relacionar el error de seguimiento con la referencia, las perturbaciones y el ruido a través de las funciones de sensitividad.
  \item Calcular el error en estado estacionario ante escalón, rampa y parábola con el tipo del sistema y las constantes $\kp$, $\kv$ y $\ka$, después de verificar la estabilidad.
  \item Tratar la retroalimentación no unitaria con el lazo abierto $\Gla$ y distinguir el error de seguimiento del error de actuación.
  \item Calcular el error que deja una perturbación, el error en espacio de estado, y reconocer cuándo una especificación de error choca con una de respuesta transitoria.
\end{objetivos}

%% ---------------------------------------------------------------------
\section{Por qué retroalimentar: sensitividad y error}
\label{sec:er-sensitividad}
%% ---------------------------------------------------------------------

Analizar un sistema de control es evaluar tres cosas: la estabilidad (sin ella las otras dos no tienen sentido), la respuesta transitoria y el \textbf{error en estado estacionario}. Las dos primeras ocuparon los capítulos~\ref{cap:po}, \ref{cap:so} y~\ref{cap:es}; este capítulo trata la tercera.

En malla abierta nada corrige una perturbación ni un cambio en los parámetros del proceso. La retroalimentación sí: reduce la sensibilidad a variaciones de la planta, mejora el rechazo de perturbaciones, reduce el error en estado estacionario y facilita ajustar el transitorio. A cambio, hace llegar al lazo el ruido del sensor. La \cref{fig:er-rdn} muestra las tres señales exógenas: la referencia $R$, la perturbación $D$ y el ruido de medición $N$. Definimos el \clave{error de seguimiento} como
\begin{equation}
  E(s)=R(s)-Y(s),
  \label{eq:er-error}
\end{equation}
lo que nos interesa del desempeño. Con $\Hs=1$, por superposición,
\begin{equation}
  E=\underbrace{\frac{1}{1+\Gc\Gp}}_{\Sens(s)}R-\frac{\Gp}{1+\Gc\Gp}D+\underbrace{\frac{\Gc\Gp}{1+\Gc\Gp}}_{\Tlc(s)}N
  =\Sens R-\Sens\Gp D+\Tlc N .
  \label{eq:er-sensitividad}
\end{equation}

\begin{figure}[ht]
\centering
\includegraphics[width=0.8\linewidth]{cap06/lazo_rdn}
\caption{Lazo con referencia $R$, perturbación $D$ a la entrada del proceso y ruido de medición $N$.}
\label{fig:er-rdn}
\end{figure}

\begin{definicion}[label=def:er-sensitividad]{Sensitividad y sensitividad complementaria}
Con retroalimentación unitaria, la \textbf{sensitividad} es $\Sens(s)=\dfrac{1}{1+\Gc\Gp}$ y la \textbf{sensitividad complementaria} es $\Tlc(s)=\dfrac{\Gc\Gp}{1+\Gc\Gp}$, que es la propia función de lazo cerrado $Y/R$. Para todo $s$,
\[
  \Sens(s)+\Tlc(s)=1 .
\]
\end{definicion}

Todos los términos de \eqref{eq:er-sensitividad} tienen el mismo denominador: la ecuación característica $1+\Gc\Gp=0$ es la misma para cualquier entrada. Y la identidad $\Sens+\Tlc=1$ dice que no se pueden hacer pequeñas las dos a la misma frecuencia. El compromiso se resuelve separando frecuencias: $\Gc\Gp$ grande donde viven las perturbaciones y las referencias (bajas frecuencias), para que $\Sens$ sea pequeña; $\Gc\Gp$ pequeña donde vive el ruido (altas frecuencias), para que $\Tlc$ lo sea.

\begin{ejemplo}[label=ej:er-sens]{Dos ganancias}
Con $\Gp=1/((s+1)(s+10))$ y $\Gc=K$, compare $K=10$ y $K=200$ ante una perturbación $d=0.5\sin(0.2t)$ y un ruido $n=0.01\sin(50t)$, con referencia constante.
\solucion La amplitud de cada aporte es la amplitud de la señal por la magnitud de su función a esa frecuencia. Con $K=200$: perturbación $0.5\,|\Sens\Gp|(0.2)=0.0024$ y ruido $0.01\,|\Tlc|(50)=0.00085$. Con $K=10$: $0.025$ y $0.00004$. La ganancia alta rechaza diez veces mejor la perturbación y deja pasar veinte veces más ruido (\cref{fig:er-sens}); aquí gana, porque el ruido es pequeño. Si su amplitud pasara de $0.28$ (sumando los aportes como cota), convendría $K=10$.
\end{ejemplo}

\begin{figure}[ht]
\centering
\includegraphics[width=0.85\linewidth]{cap06/f_sensibilidad}
\caption{Ejemplo~\ref{ej:er-sens}: $|\Sens|$ y $|\Tlc|$ para $K=10$ y $K=200$. Donde una es pequeña, la otra vale casi $1$.}
\label{fig:er-sens}
\end{figure}

\codigorepo{cap06_error/C6_Ej6_Sensibilidad.m}{aportes}{ejemplo~\ref{ej:er-sens}}

La sensitividad y la robustez se estudian a fondo en cursos posteriores. Aquí nos quedamos con su consecuencia más directa: el error que queda cuando todo se asienta. En lo que sigue (salvo en la \cref{sec:er-perturbacion}) el lazo no tiene ni perturbación ni ruido.

%% ---------------------------------------------------------------------
\section{Señales de prueba}
\label{sec:er-senales}
%% ---------------------------------------------------------------------

El \clave{error en estado estacionario} es el límite del error de seguimiento cuando $t\to\infty$, ante una señal de prueba. Las tres tradicionales son (\cref{fig:er-senales}):
\begin{itemize}
  \item \textbf{Escalón}, $r(t)=\escalon(t)$, $R(s)=1/s$: mantener una posición o un valor constante (un satélite geoestacionario, una temperatura).
  \item \textbf{Rampa}, $r(t)=t\,\escalon(t)$, $R(s)=1/s^2$: seguir un blanco que se mueve a velocidad constante (una antena que sigue un satélite).
  \item \textbf{Parábola}, $r(t)=\frac{t^2}{2}\escalon(t)$, $R(s)=1/s^3$: seguir un blanco con aceleración constante.
\end{itemize}

\begin{figure}[ht]
\centering
\includegraphics[width=\linewidth]{cap06/f_senales}
\caption{Las tres señales de prueba.}
\label{fig:er-senales}
\end{figure}

\begin{cuidado}[Primero la estabilidad]
El teorema del valor final solo vale si $sE(s)$ tiene todos sus polos en el semiplano izquierdo, es decir, si el lazo cerrado es estable. Calcular un ``error en estado estacionario'' de un lazo inestable da un número que no significa nada. \alerta{Verifique la estabilidad primero} (\cref{cap:es}).
\end{cuidado}

%% ---------------------------------------------------------------------
\section{Tipo del sistema y constantes de error}
\label{sec:er-tipo}
%% ---------------------------------------------------------------------

Con retroalimentación unitaria (\cref{fig:er-unitario}), $\Tlc=G/(1+G)$ y $E(s)=\dfrac{1}{1+G(s)}R(s)$. Si el lazo cerrado es estable, el teorema del valor final da
\begin{equation}
  \boxed{\;\ess=\lim_{t\to\infty}e(t)=\lim_{s\to0}\,s\,\frac{R(s)}{1+G(s)}\;}
  \label{eq:er-ess}
\end{equation}
Aquí $G(s)$ es el \textbf{lazo abierto}: la función que, con retroalimentación unitaria, va del error a la salida (\cref{def:rf-malla}; si la retroalimentación no es unitaria, ver la \cref{sec:er-no-unitaria}).

\begin{figure}[ht]
\centering
\includegraphics[width=0.55\linewidth]{cap06/lazo_unitario}
\caption{Lazo con retroalimentación unitaria.}
\label{fig:er-unitario}
\end{figure}

Escribamos el lazo abierto en forma de Bode, separando los polos en el origen:
\begin{equation}
  G(s)=\frac{K\,(T_1s+1)(T_2s+1)\cdots(T_ms+1)}{s^N(\tau_1s+1)(\tau_2s+1)\cdots(\tau_{n-N}s+1)} .
  \label{eq:er-bode}
\end{equation}

\begin{definicion}[label=def:er-tipo]{Tipo del sistema}
El \textbf{tipo} es el número $N$ de polos en el origen del lazo abierto, es decir, el número de integradores del lazo. $K$ es su ganancia de Bode: el valor de $s^NG(s)$ en $s=0$.
\end{definicion}

Cada señal de prueba define una \textbf{constante estática de error}; cuanto mayor, menor el error:
\begin{itemize}
  \item \textbf{Escalón}, $R=1/s$: $\ess=\dfrac{1}{1+\lim_{s\to0}G(s)}=\dfrac{1}{1+\kp}$, con la \clave{constante de posición} $\kp=\lim_{s\to0}G(s)$.
  \item \textbf{Rampa}, $R=1/s^2$: $\ess=\lim_{s\to0}\dfrac{1}{s+sG(s)}=\dfrac{1}{\kv}$, con la \clave{constante de velocidad} $\kv=\lim_{s\to0}sG(s)$.
  \item \textbf{Parábola}, $R=1/s^3$: $\ess=\dfrac{1}{\ka}$, con la \clave{constante de aceleración} $\ka=\lim_{s\to0}s^2G(s)$.
\end{itemize}
Con \eqref{eq:er-bode}, $s^iG(s)\to K/s^{N-i}$: la constante vale $K$ si $N=i$, infinito si $N>i$ y cero si $N<i$. Eso da la tabla que resume el capítulo.

\begin{table}[H]
\centering
\caption{Error en estado estacionario según el tipo del lazo abierto (lazo cerrado estable, retroalimentación unitaria). Para una entrada de amplitud $A$ (escalón $A$, rampa $At$, parábola $At^2/2$), el error se multiplica por $A$.}
\label{tab:er-tabla}
\renewcommand{\arraystretch}{1.5}
\begin{tabular}{@{}lcccccc@{}}
\toprule
 & $\kp$ & $\kv$ & $\ka$ & escalón & rampa & parábola\\
\midrule
Tipo 0 & $K$ & $0$ & $0$ & $\dfrac{1}{1+\kp}$ & $\infty$ & $\infty$\\
Tipo 1 & $\infty$ & $K$ & $0$ & $0$ & $\dfrac{1}{\kv}$ & $\infty$\\
Tipo 2 & $\infty$ & $\infty$ & $K$ & $0$ & $0$ & $\dfrac{1}{\ka}$\\
\bottomrule
\end{tabular}
\end{table}

\begin{ideaclave}
\textbf{Cada integrador vale una derivada de la referencia.}\enspace Un lazo de tipo $N$ sigue sin error una referencia polinomial de grado menor que $N$, con error finito una de grado $N$, y pierde una de grado mayor. El error se lee en la tabla con solo dos datos: el tipo y una ganancia.
\end{ideaclave}

\begin{figure}[ht]
\centering
\includegraphics[width=\linewidth]{cap06/f_tipos}
\caption{Lazos abiertos de tipo 0, 1 y 2 ($\frac{4}{(s+1)(s+2)}$, $\frac{4}{s(s+2)}$ y $\frac{4(s+0.5)}{s^2(s+2)}$) ante escalón y rampa unitarios. Errores: $1/3$ y $\infty$; $0$ y $0.5$; $0$ y $0$.}
\label{fig:er-tipos}
\end{figure}

\begin{ejemplo}[label=ej:er-motor]{Un motor con retroalimentación de velocidad}
En el servo de la \cref{fig:er-motor}~\citep{DoBi08}, halle $\ess$ ante una rampa unitaria en función de $K$, $K_b$ y $K_m$, y escoja $K$ para $\ess=1$ con $K_m=10$ y $K_b=0.05$.
\solucion El lazo interno da $\Omega/(KE)=K_m/(s+0.01+K_bK_m)$, y con el integrador,
\[
  G(s)=\frac{KK_m}{s(s+0.01+K_bK_m)} :
\]
tipo 1, con $\kv=\dfrac{KK_m}{0.01+K_bK_m}$ y $\ess=\dfrac{0.01+K_bK_m}{KK_m}$. Con los datos, $\ess=0.51/(10K)=1$ da $K=0.051$. (El lazo cerrado es de segundo orden con coeficientes positivos: estable para todo $K>0$.)
\end{ejemplo}

\begin{figure}[ht]
\centering
\includegraphics[width=0.85\linewidth]{cap06/lazo_motor_kb}
\caption{Ejemplo~\ref{ej:er-motor}: servo de posición con retroalimentación interna de velocidad.}
\label{fig:er-motor}
\end{figure}

\codigorepo{cap06_error/C6_Ej1_TipoError.m}{tipo}{tipo, constantes y estabilidad de cuatro lazos con la función \cmd{constantesError}}
\codigorepo{cap06_error/C6_Ej1_TipoError.m}{simulacion}{el error de rampa simulado}

\begin{pausa}
Un lazo abierto $G=\dfrac{5(s+1)}{s^3(s+2)}$ tiene $\kp=\kv=\ka=\infty$. ¿Sigue sin error una parábola?
\respuesta{No se puede saber con la tabla: el polinomio característico es $s^4+2s^3+5s+5$, al que le falta el término en $s^2$, así que el lazo cerrado es inestable. La salida no sigue nada.}
\end{pausa}

%% ---------------------------------------------------------------------
\section{Retroalimentación no unitaria}
\label{sec:er-no-unitaria}
%% ---------------------------------------------------------------------

Si el sensor tiene dinámica o ganancia, $\Hs(s)\neq1$, y al comparador ya no llega el error de seguimiento, sino el \clave{error de actuación}
\[
  E_a(s)=R(s)-\Hs(s)Y(s),
\]
la señal que entra a $G$. Aplicar la \cref{tab:er-tabla} a la malla abierta $\Gma=G\Hs$ da el error de \emph{actuación} en estado estacionario, no el de seguimiento. Para el de seguimiento, el lazo se reescribe con retroalimentación unitaria (ejemplo~\ref{ej:so-equivalente}) con el \textbf{lazo abierto}
\begin{equation}
  \Gla(s)=\frac{G(s)}{1+G(s)\Hs(s)-G(s)},
  \label{eq:er-gla}
\end{equation}
y con él se usan el tipo, las constantes y la tabla. Si $\Hs=1$, $\Gla=G$.

\begin{ejemplo}[label=ej:er-gh]{Dos errores que no coinciden}
En la \cref{fig:er-gh}, $G=\dfrac{100}{s+1}$ y $\Hs=\dfrac{1}{s+5}$. Halle el error de actuación y el de seguimiento ante un escalón unitario.
\solucion El lazo cerrado es $\Tlc=\dfrac{G}{1+G\Hs}=\dfrac{100(s+5)}{s^2+6s+105}$, estable (polos en $-3\pm9.8j$), con $y_\infty=\Tlc(0)=500/105=4.76$.

\emph{Con la malla abierta}: $G\Hs=\dfrac{100}{(s+1)(s+5)}$, tipo 0, $\kp=20$, y $e_{a,\infty}=1/21=0.048$. Es correcto, pero es el error de actuación: en estado estacionario $\Hs(0)y_\infty=4.76/5=0.952\approx r$.

\emph{Con el lazo abierto}: \eqref{eq:er-gla} da
\[
  \Gla(s)=\frac{100(s+5)}{s^2-94s-395},\qquad \kp=\Gla(0)=-\frac{500}{395}=-1.266,
\]
y $\ess=\dfrac{1}{1+\kp}=\dfrac{1}{1-1.266}=-3.76=1-y_\infty$. La salida termina en $4.76$ veces la referencia (\cref{fig:er-no-unitaria}). El signo negativo dice que la salida \emph{sobrepasa} la referencia.

$\Gla$ tiene un polo en $+98$: no es un sistema físico que haya que estabilizar, sino una función auxiliar que reproduce el mismo lazo cerrado; el lazo cerrado sigue siendo estable. La causa del error es la ganancia DC del sensor, $\Hs(0)=1/5$: el lazo lleva $\Hs y$ a la referencia, no $y$. Si el sensor se recalibra como $\Hs=k/(s+5)$, el error de seguimiento se anula con $\Tlc(0)=500/(5+100k)=1$, es decir, $k=4.95$.
\end{ejemplo}

\begin{figure}[ht]
\centering
\includegraphics[width=0.55\linewidth]{cap06/lazo_gh}
\caption{Ejemplo~\ref{ej:er-gh}: lazo con un sensor dinámico.}
\label{fig:er-gh}
\end{figure}

\begin{figure}[ht]
\centering
\includegraphics[width=0.85\linewidth]{cap06/f_no_unitaria}
\caption{Ejemplo~\ref{ej:er-gh}: el error de actuación tiende a $1/21$ y el de seguimiento a $-3.76$.}
\label{fig:er-no-unitaria}
\end{figure}

\codigorepo{cap06_error/C6_Ej2_NoUnitaria.m}{lazoabierto}{malla abierta, lazo abierto y los dos errores}
\codigorepo{cap06_error/C6_Ej2_NoUnitaria.m}{simulacion}{la simulación y el sensor recalibrado}

\begin{cuidado}[Malla abierta no es lazo abierto]
La estabilidad y el diagrama de Bode se estudian con la malla abierta $G\Hs$; el error de seguimiento, con el lazo abierto $\Gla$. Coinciden solo si $\Hs=1$. Si $\Hs(0)=1$ (un sensor bien calibrado), el error ante escalón también coincide, aunque las funciones sean distintas.
\end{cuidado}

%% ---------------------------------------------------------------------
\section{Error debido a perturbaciones}
\label{sec:er-perturbacion}
%% ---------------------------------------------------------------------

Con $R=0$ y $N=0$, \eqref{eq:er-sensitividad} da el error que deja una perturbación a la entrada del proceso:
\begin{equation}
  E(s)=-\frac{\Gp(s)}{1+\Gc(s)\Gp(s)}\,D(s).
  \label{eq:er-perturbacion}
\end{equation}
Si la perturbación entra a la salida del proceso, $E=-\dfrac{1}{1+\Gc\Gp}D$. En los dos casos, el error ante una perturbación escalón se anula si $\Gc$ tiene un integrador: \alerta{importa dónde está el integrador}. Uno en el controlador, \emph{antes} del punto donde entra la perturbación, la rechaza; uno en la planta, después de ese punto, no.

\begin{ejemplo}[label=ej:er-perturbacion]{Proporcional contra PI}
Con $\Gp=\dfrac{1}{(s+1)(s+2)}$ y una perturbación escalón de $0.5$, halle el error con $\Gc=K$ y con $\Gc=K+K_i/s$.
\solucion \emph{Proporcional, perturbación a la entrada}: $\ess=-0.5\,\dfrac{\Gp(0)}{1+K\Gp(0)}=-\dfrac{0.5}{2+K}$: $-0.167$, $-0.042$ y $-0.0049$ para $K=1$, $10$ y $100$. Baja con $K$ pero no se anula. \emph{A la salida}: $\ess=-\dfrac{0.5}{1+K/2}=-\dfrac{1}{2+K}$, el doble, porque la perturbación no pasa por $\Gp(0)=1/2$.

\emph{PI}: $\ess=-0.5\lim_{s\to0}\dfrac{\Gp}{1+(K+K_i/s)\Gp}=-0.5\lim_{s\to0}\dfrac{s\Gp}{s+(Ks+K_i)\Gp}=0$ para cualquier $K_i\neq0$ que deje el lazo estable. Con $K=10$, la ecuación característica es $s^3+3s^2+12s+K_i=0$ y Routh exige $0<K_i<36$ (\cref{fig:er-perturbacion}).
\end{ejemplo}

\begin{figure}[ht]
\centering
\includegraphics[width=\linewidth]{cap06/f_perturbacion}
\caption{Ejemplo~\ref{ej:er-perturbacion}: error por una perturbación escalón de $0.5$.}
\label{fig:er-perturbacion}
\end{figure}

\codigorepo{cap06_error/C6_Ej3_Perturbacion.m}{proporcional}{el controlador proporcional, perturbación a la entrada y a la salida}
\codigorepo{cap06_error/C6_Ej3_Perturbacion.m}{pi}{el controlador PI}

\begin{modelo}[Modelo de Simulink: \texttt{crearModelo\_Perturbacion}]
\texttt{codigo/cap06\_error/crearModelo\_Perturbacion.m} arma el lazo con un PI y una perturbación escalón de $0.5$ en $t=5$\,s. La variable \texttt{alEntrada} decide si entra a la entrada o a la salida del proceso; con \texttt{Ki = 0} el controlador es proporcional.
\end{modelo}

%% ---------------------------------------------------------------------
\section{Error en espacio de estado}
\label{sec:er-estado}
%% ---------------------------------------------------------------------

Si el proceso está en espacio de estado (\cref{cap:ee}), el error en estado estacionario se obtiene sin pasar por la función de transferencia: con el lazo cerrado y una referencia constante, el estado estacionario cumple $\dot x=0$.

\begin{ejemplo}[label=ej:er-estado]{Estado estacionario directo}
Un proceso con $\Am=\begin{bmatrix}0&1\\-2&-3\end{bmatrix}$, $\Bm=[0\;\,1]^\top$, $\Cm=[2\;\,0]$ se controla con $u=r-y$. Halle $\ess$ ante un escalón unitario. Después agregue un estado integrador $\dot x_3=r-y$ con $u=r-y+K_ix_3$.
\solucion Con la función de transferencia: $G=\dfrac{2}{s^2+3s+2}$, tipo 0, $\kp=1$, $\ess=0.5$. En espacio de estado: $\dot x=(\Am-\Bm\Cm)x+\Bm r$ y, con $\dot x=0$,
\[
  x_{ss}=-(\Am-\Bm\Cm)^{-1}\Bm\,r=\begin{bmatrix}0.25\\0\end{bmatrix},\qquad y_{ss}=\Cm x_{ss}=0.5,\qquad \ess=0.5 .
\]
Con el estado integrador, el sistema aumentado es
\[
  \begin{bmatrix}\dot x\\\dot x_3\end{bmatrix}=
  \begin{bmatrix}\Am-\Bm\Cm&K_i\Bm\\-\Cm&0\end{bmatrix}\begin{bmatrix}x\\x_3\end{bmatrix}+\begin{bmatrix}\Bm\\1\end{bmatrix}r,
\]
con polinomio característico $s^3+3s^2+4s+2K_i$: estable para $0<K_i<6$. En estado estacionario $\dot x_3=r-y=0$: el error se anula para todo $K_i$ que deje el lazo estable. Con $K_i=3$, los polos son $-2.38$ y $-0.31\pm1.56j$.
\end{ejemplo}

\codigorepo{cap06_error/C6_Ej4_Estado.m}{estacionario}{el estado estacionario con $\dot x=0$}
\codigorepo{cap06_error/C6_Ej4_Estado.m}{integral}{el estado integrador}

El estado integrador convierte el lazo en uno de tipo 1. Es la misma idea de la acción integral de un PI, y la que se usa con la retroalimentación de estado en la \cref{sec:re-integral}.

%% ---------------------------------------------------------------------
\section{Especificaciones de error y de transitorio}
\label{sec:er-diseno}
%% ---------------------------------------------------------------------

Una especificación de error se traduce en una cota sobre una constante: ``$\ess\le0.05$ ante rampa unitaria'' es $\kv\ge20$. Con un controlador proporcional, la única perilla es $K$, y la misma perilla mueve los polos del lazo cerrado. Las dos especificaciones pueden chocar.

\begin{ejemplo}[label=ej:er-conflicto]{Un servo que no puede con las dos}
Con $G_p=\dfrac{1}{s(s+4)(s+8)}$ y $\Gc=K$, se pide $\ess\le0.05$ ante rampa unitaria y $\OS\le20$\,\% ante escalón.
\solucion Tipo 1 con $\kv=K/32$: el error pide $K\ge640$. Pero la ecuación característica $s^3+12s^2+32s+K=0$ exige $K<12\cdot32=384$: \alerta{con $K$ solo, ni siquiera es estable}. Y el sobrepico pide un $K$ mucho menor: para $\zeta=0.49$ en los polos dominantes ($17$\,\% estimado, con margen; la simulación da $16$\,\%), $K=67.8$, que da $\kv=2.1$ y $\ess=0.47$ (\cref{fig:er-atraso}, azul).

Un \textbf{compensador de atraso} $\Gc=K\dfrac{s+a}{s+b}$ con $a>b>0$ multiplica $\kv$ por $a/b$. Si el par polo-cero está muy cerca del origen, en la zona de los polos dominantes $(s+a)/(s+b)\approx1$, y el transitorio casi no cambia. Se escoge $K=67.8$ por el transitorio, y $a/b=1.1\cdot640/67.8=10.4$ (con un $10$\,\% de margen); con $b=0.005$, $a=0.0519$. El lazo compensado tiene $\kv=22$ y $\ess=0.045$, con $\OS=19$\,\%; el precio es una cola lenta (un polo en $-0.053$ casi cancelado por el cero en $-0.052$): el tiempo de establecimiento sube de $3.1$ a $4.8$\,s, y el error de rampa tarda unos $116$\,s en quedar a menos del $2$\,\% de su valor final.
\end{ejemplo}

\begin{figure}[ht]
\centering
\includegraphics[width=\linewidth]{cap06/f_atraso}
\caption{Ejemplo~\ref{ej:er-conflicto}: con solo $K=67.8$ (azul) y con el atraso (coral). El atraso deja casi igual la respuesta al escalón y baja el error de rampa de $0.47$ a $0.045$, con una cola lenta.}
\label{fig:er-atraso}
\end{figure}

\codigorepo{cap06_error/C6_Ej5_Atraso.m}{conflicto}{el conflicto entre las dos especificaciones}
\codigorepo{cap06_error/C6_Ej5_Atraso.m}{atraso}{el compensador de atraso}

El diseño de atrasos en el dominio de la frecuencia es tema de la \cref{sec:rf-atraso}; el PI, que lleva $b$ hasta cero, del \cref{cap:pn}.

%% ---------------------------------------------------------------------
\section{Para llevar}
\label{sec:er-cierre}
%% ---------------------------------------------------------------------

\begin{parallevar}
\begin{enumerate}
  \item $E=\Sens R-\Sens\Gp D+\Tlc N$ con $\Sens+\Tlc=1$: lazo abierto grande donde están la referencia y las perturbaciones, pequeño donde está el ruido.
  \item Primero la estabilidad. Después, $\ess=\lim_{s\to0}sR(s)/(1+G(s))$.
  \item El tipo es el número de integradores del lazo abierto. $\kp=\lim G$, $\kv=\lim sG$, $\ka=\lim s^2G$; escalón $1/(1+\kp)$, rampa $1/\kv$, parábola $1/\ka$.
  \item Con retroalimentación no unitaria, el error de seguimiento se calcula con el lazo abierto $\Gla=G/(1+G\Hs-G)$; la malla abierta $G\Hs$ da el error de actuación.
  \item Una perturbación escalón se rechaza por completo si hay un integrador \emph{antes} del punto donde entra.
  \item En espacio de estado, el estado estacionario sale de $\dot x=0$; un estado integrador del error vuelve el lazo de tipo 1.
  \item Error y transitorio pueden chocar con un solo parámetro; un atraso sube $\kv$ por el factor $a/b$ a costa de una cola lenta.
\end{enumerate}
\end{parallevar}

\begin{profundiza}[Correcciones respecto a la versión anterior del capítulo]
(i) En las secciones de $\kv$ y $\ka$ se escribía $\kv=\lim_{s\to0}G(s)$ y $\ka=\lim_{s\to0}G(s)$ antes del desarrollo; son $\lim sG$ y $\lim s^2G$. (ii) La sensitividad complementaria se llamaba $C(s)$; se usa $\Tlc(s)$, que es además la función de lazo cerrado, para no confundirla con $\Cm$. (iii) El error con retroalimentación no unitaria se trata ahora con el lazo abierto $\Gla$ del \cref{cap:so} y la distinción entre error de actuación y de seguimiento, que la versión anterior dejaba para un ejercicio. (iv) La tabla resumen tenía una columna vacía de más y no incluía las constantes; se completó, y se dejó explícito que solo vale con el lazo cerrado estable. (v) Los ejemplos físicos de las constantes (``la temperatura de salida'') se retiraron por imprecisos, y el código de Mathematica de las señales de prueba se reemplazó por una figura generada con el resto. (vi) Se resolvió el ejemplo del lazo con sensor dinámico, que estaba solo enunciado, y se agregaron el error por perturbaciones, el error en espacio de estado y el conflicto entre especificaciones.
\end{profundiza}

%% ---------------------------------------------------------------------
\section{Problemas}
\label{sec:er-problemas}
%% ---------------------------------------------------------------------

Los problemas marcados con \etiquetaM{} requieren \matlab/\simulink; los marcados con \etiquetaEstrella{} son de profundización.

\begin{problema}[tipo y constantes]\label{prob:er-tipos}
Con retroalimentación unitaria y lazo abierto (i) $\dfrac{20(s+2)}{(s+1)(s+5)}$, (ii) $\dfrac{50}{s(s+3)(s+10)}$, (iii) $\dfrac{100(s+4)(s+6)}{s^2(s+5)(s+8)}$ y (iv) $\dfrac{5(s+1)}{s^3(s+2)}$: (a) verifique la estabilidad con Routh; (b) halle el tipo y $\kp$, $\kv$, $\ka$; (c) calcule $\ess$ ante $r=5$, $r=2t$ y $r=\frac32t^2$ y tabule; (d) uno de los cuatro no admite el análisis: ¿cuál y por qué? (e) \etiquetaM{} verifique por simulación los casos admisibles.
\end{problema}

\begin{problema}[error contra estabilidad]\label{prob:er-posicion}
Un servo tiene lazo abierto $\dfrac{K}{s(s+4)(s+8)}$ con retroalimentación unitaria. (a) Halle el tipo y $\kv(K)$. (b) Halle el menor $K$ con $\ess\le0.1$ ante rampa unitaria. (c) Halle el rango de $K$ estable. (d) ¿Hay valores que cumplan las dos? (e) Para el $K$ de (b), halle los polos, $\zeta$ y $\omega_n$ del par dominante, y estime $\OS$ y $\ts$; \etiquetaM{} verifique con \cmd{stepinfo}. (f) ¿Por qué no es aceptable reducir el error solo con la ganancia?
\end{problema}

\begin{problema}[perturbaciones]\label{prob:er-perturbacion}
En la \cref{fig:er-rdn}, con $\Hs=1$, $N=R=0$, $\Gp=\dfrac{1}{(s+1)(s+2)}$ y una perturbación escalón de $0.5$ a la entrada del proceso: (a) halle $\ess$ con $\Gc=K$ para $K=1$, $10$ y $100$; (b) muestre que con $\Gc=K+K_i/s$ el error es nulo para todo $K_i$ estable, y diga qué propiedad de $\Gc$ lo produce; (c) halle el rango de $K_i$ estable con $K=10$; (d) repita (a) con la perturbación a la salida y explique la diferencia; (e) \etiquetaM{} verifique con \texttt{crearModelo\_Perturbacion}.
\end{problema}

\begin{problema}[sensor dinámico]\label{prob:er-gh}
Con $G=\dfrac{100}{s+1}$ y $\Hs=\dfrac{1}{s+5}$: (a) verifique la estabilidad; (b) halle $\kp$ de la malla abierta y el error de actuación ante escalón; (c) halle $y_\infty$ y el error de seguimiento; (d) explique por qué difieren y cuándo coinciden; (e) \etiquetaM{} simule $r$, $y$, $e_a$ y $e$; ¿cuál error da la \cref{tab:er-tabla} con la malla abierta?
\end{problema}

\begin{problema}[lazo abierto equivalente]\label{prob:er-gla}
Con los datos del problema~\ref{prob:er-gh}: (a) muestre que el lazo con $G$ y $\Hs$ tiene el mismo $\Tlc$ que uno unitario con $\Gla$ de \eqref{eq:er-gla}; (b) calcule $\Gla$ y sus polos; (c) explique cómo puede tener un polo en el semiplano derecho si el lazo cerrado es estable; (d) calcule $\kp$ de $\Gla$ y $\ess=1/(1+\kp)$, y compárelo con (c) del problema anterior; (e) interprete el signo; (f) con $\Hs=k/(s+5)$, halle el $k$ que anula el error y discuta la importancia de calibrar el sensor.
\end{problema}

\begin{problema}[espacio de estado]\label{prob:er-estado}
Con $\Am$, $\Bm$, $\Cm$ del ejemplo~\ref{ej:er-estado} y $u=r-y$: (a) halle $G$, el tipo, $\kp$ y $\ess$; (b) obtenga lo mismo con $\dot x=0$; (c) escriba el sistema aumentado con $\dot x_3=r-y$ y $u=r-y+K_ix_3$; (d) halle el rango de $K_i$ estable; (e) muestre que el error es nulo y relaciónelo con el tipo; (f) \etiquetaM{} simule con $K_i$ en la mitad del rango.
\end{problema}

\begin{problema}[atraso]\label{prob:er-atraso}
Con $G_p=\dfrac{1}{s(s+4)(s+8)}$ se pide $\ess\le0.05$ ante rampa unitaria y $\OS\le20$\,\%. (a) Muestre que ningún $\Gc=K$ cumple las dos y cuantifique el conflicto. (b) Muestre que $\Gc=K\frac{s+a}{s+b}$ multiplica $\kv$ por $a/b$ y casi no mueve los polos dominantes si $a$ y $b$ están cerca del origen. (c) Escoja $K$ por el transitorio y $(a,b)$ por el error. (d) Verifique la estabilidad con Routh sobre el polinomio de cuarto orden. (e) \etiquetaM{} Simule escalón y rampa, verifique las especificaciones y cuantifique la cola lenta. (f) ¿Qué precio se paga y cuándo sería inaceptable?
\end{problema}

\begin{problema}[sensitividad]\label{prob:er-sens}
Con $\Gp=\dfrac{1}{(s+1)(s+10)}$ y $\Gc=K$: (a) muestre que $\Sens+\Tlc=1$; (b) \etiquetaM{} tabule $|\Sens|$ y $|\Tlc|$ en $\omega=0.1$, $1$, $10$ y $100$\,rad/s para $K=10$ y $K=200$; (c) halle la frecuencia donde $|\Sens|=|\Tlc|$ y relaciónela con el cruce de ganancia; (d) con $d=0.5\sin(0.2t)$ a la entrada y $n=0.01\sin(50t)$, halle la amplitud de cada aporte al error; (e) ¿cuál $K$ da menor error y por qué? (f) \etiquetaEstrella{} ¿Desde qué amplitud de ruido conviene $K=10$?
\end{problema}

%% ---------------------------------------------------------------------
\section{Soluciones seleccionadas}
%% ---------------------------------------------------------------------

\paragraph{Problema~\ref{prob:er-tipos}}
(a) Polinomios $s^2+26s+45$, $s^3+13s^2+30s+50$, $s^4+13s^3+140s^2+1000s+2400$ (estables) y $s^4+2s^3+5s+5$ (inestable: falta $s^2$). (b) (i) tipo 0, $\kp=8$; (ii) tipo 1, $\kv=5/3$; (iii) tipo 2, $\ka=60$. (c) (i) $5/9$, $\infty$, $\infty$; (ii) $0$, $1.2$, $\infty$; (iii) $0$, $0$, $0.05$. (d) El (iv): la tabla daría error nulo para todo, pero la salida diverge.

\paragraph{Problema~\ref{prob:er-posicion}}
(a) Tipo 1, $\kv=K/32$. (b) $K=320$. (c) $0<K<384$. (d) $320\le K<384$. (e) Polos $-11.6$ y $-0.19\pm5.25j$: $\zeta=0.037$, $\omega_n=5.25$, $\OS\approx89$\,\% estimado ($81$\,\% simulado) y $\ts\approx21$\,s. (f) El lazo queda casi en el límite de estabilidad.

\paragraph{Problema~\ref{prob:er-perturbacion}}
(a) $-0.167$, $-0.042$, $-0.0049$. (b) El integrador de $\Gc$, que está antes de la perturbación. (c) $0<K_i<36$. (d) $-0.333$, $-0.083$, $-0.0098$: a la salida, la perturbación no se atenúa por $\Gp(0)=1/2$.

\paragraph{Problema~\ref{prob:er-gh}}
(a) Polos $-3\pm9.8j$. (b) $\kp=20$, $e_a=1/21$. (c) $y_\infty=100/21=4.76$, $\ess=-3.76$. (d) Coinciden si $\Hs(0)=1$.

\paragraph{Problema~\ref{prob:er-gla}}
(b) $\Gla=\dfrac{100(s+5)}{s^2-94s-395}$, polos $98.0$ y $-4.03$. (d) $\kp=-1.266$ y $\ess=-3.76$. (f) $k=4.95$.

\paragraph{Problema~\ref{prob:er-estado}}
(a) $G=\frac{2}{s^2+3s+2}$, tipo 0, $\kp=1$, $\ess=0.5$. (d) $s^3+3s^2+4s+2K_i$: $0<K_i<6$. (e) $\dot x_3=r-y=0$ en estado estacionario; el lazo es de tipo 1.

\paragraph{Problema~\ref{prob:er-atraso}}
(a) $\ess\le0.05$ pide $K\ge640$; la estabilidad, $K<384$; $\OS\le20$\,\%, $K\lesssim74$. (c) Con $K=67.8$, $b=0.005$ y $a=0.0519$: $\kv=22$. (e) $\OS=19$\,\%, $\ess=0.045$; $\ts$ pasa de $3.1$ a $4.8$\,s y el error de rampa tarda unos $116$\,s en asentarse.

\paragraph{Problema~\ref{prob:er-sens}}
(b) $K=10$: $|\Sens|=0.50$, $0.65$, $1.04$, $1.00$ y $|\Tlc|=0.50$, $0.46$, $0.07$, $0.001$; $K=200$: $|\Sens|=0.048$, $0.068$, $0.91$, $1.02$ y $|\Tlc|=0.95$, $0.96$, $1.29$, $0.020$. (c) Con $K=200$, en $12.5$\,rad/s, el cruce de $|\Gc\Gp|=1$; con $K=10$, $|\Gc\Gp|\le1$ en todas las frecuencias y las dos solo se igualan en $\omega\to0$. (d) $K=10$: $0.025$ y $0.00004$; $K=200$: $0.0024$ y $0.00085$. (f) Ruido de amplitud mayor que $0.28$.

\respuestaspausas

